\chapter{1943:George de Hevesy}

\label{chap:chemistry1943}

The Nobel Prize in Chemistry for the year 1943 was awarded to George de Hevesy.

\textbf{Motivation:}

for his work on the use of isotopes as tracers in the study of chemical processes

\section{George de Hevesy}

\subsection*{Early Life and Education}

George Charles de Hevesy, born György Hevesy, was born on 1885-08-01 in Budapest, Hungary, then part of Austria-Hungary.
He was a Hungarian-Swedish radiochemist awarded the Nobel Prize in Chemistry in 1943 for his work on the use of isotopes as tracers in the study of chemical processes.

\subsection*{Career and Major Research}

de Hevesy studied chemistry in Budapest, Freiburg, Berlin, and Zürich, and received his doctorate at the University of Freiburg in 1908.
He worked with Ernest Rutherford in Manchester from 1911 to 1913, where his unsuccessful attempts to separate the radioactive isotope radium D from lead led him to realize that isotopes are chemically inseparable and could instead be used to trace the path of an element through chemical and biological processes.
In 1923, together with Dirk Coster, he discovered the element hafnium (atomic number 72) in zirconium minerals, using X-ray spectroscopy.
From 1920 he worked at the Niels Bohr Institute in Copenhagen, and in 1926 he became professor at the University of Freiburg; he returned to Copenhagen in 1934.
He applied radioactive indicators as tracers in chemistry, biology, and medicine, measuring the absorption and distribution of elements in plants and in animals.
In 1943, during the German occupation of Denmark, he fled to Sweden, where he worked at the University of Stockholm and at Uppsala until his retirement.

\subsection*{Nobel Prize-winning Work}

The work for which George de Hevesy was awarded the Nobel Prize in Chemistry in 1943 was described by the Nobel Committee as: "for his work on the use of isotopes as tracers in the study of chemical processes".

de Hevesy showed that a radioactive isotope of an element behaves chemically exactly like the ordinary element and can therefore be used to follow that element through a chemical reaction, a plant, or a living organism.

\subsection*{Significance and Impact}

The tracer method developed by de Hevesy has become one of the most widely used techniques in all of science.
Radioisotopes are now used to follow the course of chemical reactions, the transport of substances in the body, and the fate of atoms in industrial processes.

\subsection*{Later Career and Honors}

de Hevesy continued his research in Sweden after the war and returned to Freiburg in 1961 after his retirement.
He received the Copley Medal of the Royal Society in 1966 and held honorary doctorates from many universities, and the George Hevesy Medal is awarded for achievements in radioanalytical chemistry.
He died on 1966-07-05 in Freiburg im Breisgau, Germany.

\subsection*{Legacy}

The legacy of George de Hevesy endures through the lasting impact of his scientific contributions.
The isotope tracer method that he invented continues to underpin chemistry, biology, and medicine.

\subsection*{Modern Developments}

de Hevesy's invention of radioactive tracers founded the field of radiochemistry and isotope labeling, now indispensable across biochemistry, medicine, and the life sciences. Tracer techniques are used in metabolic studies, drug development, PET imaging, and environmental science. de Hevesy was also the father of X-ray fluorescence analysis, a standard tool in materials science.

\subsection*{Selected Publications}

\begin{itemize}

\item Coster, D., \& Hevesy, G. (1923). "On the Missing Element of Atomic Number 72." Nature, 111, 79.

\item Hevesy, G., \& Paneth, F. (1926). "A Manual of Radioactivity." Oxford: Oxford University Press.

\end{itemize}

\clearpage

\section*{Chapter Summary}

The 1943 prize recognized George de Hevesy for his work on the use of isotopes as tracers in the study of chemical processes. His tracer method transformed biochemistry and medicine and led to the discovery of the element hafnium.
